% Part:sets-functions-relations
% Chapter: size-of-sets
% Section: equinumerous-sets
% Hindi translation (hi-Deva-IN) of Open Logic commit 9620cc73f9c8e0ad003c514a5d3748f29611c4c0.

\documentclass[../../../include/open-logic-section]{subfiles}

\begin{document}

\olfileid[hi]{sfr}{siz}{equ}

\olsection{समसंख्यता}

समुच्चयों के ``आकार'' की हमारी एक सहज धारणा है, जो परिमित समुच्चयों के लिए
ठीक काम करती है। लेकिन अनंत समुच्चयों का क्या? यदि हम \emph{किसी भी} आकार
के दो समुच्चयों के आकारों की तुलना करने का कोई विधिवत ढंग चाहते हैं, तो
पहले यह परिभाषित करना उचित है कि समुच्चय समान आकार के कब होते हैं। फ्रेगे
इसे इस प्रकार कहते हैं:
\begin{quote}
  यदि कोई परिचारक यह सुनिश्चित करना चाहता है कि उसने मेज पर ठीक उतनी ही
  छुरियाँ रखी हैं जितनी थालियाँ, तो उसे दोनों में से किसी को भी गिनने की
  आवश्यकता नहीं है; वह बस प्रत्येक थाली के दाहिने एक छुरी इस प्रकार रख दे
  कि मेज की प्रत्येक छुरी किसी थाली के दाहिने हो। इस तरह थालियाँ और छुरियाँ
  एक-दूसरे के साथ अद्वितीय रूप से सहसंबद्ध होती हैं, और वह भी उसी स्थानिक
  संबंध के माध्यम से। \citep[\S70]{Frege1884}
\end{quote}
इस अनुच्छेद की अंतर्दृष्टि को निम्न विधिवत परिभाषा से स्पष्ट किया जा सकता है:

\begin{defn}\ollabel{comparisondef}
  $A$ और $B$ \emph{समसंख्यक} हैं, जिसे $\cardeq{A}{B}$ लिखते हैं, यदि
  और केवल यदि कोई !!a{bijection} $f \colon A \to B$ हो।
\end{defn}

\begin{prop}\ollabel{equinumerosityisequi}
समसंख्यता एक तुल्यता संबंध है।
\end{prop}

\begin{proof} 
हमें दिखाना है कि समसंख्यता स्वतुल्य, सममित और संक्रामक है। मान लीजिए
$A, B$ और $C$ समुच्चय हैं।

\emph{स्वतुल्यता।} तत्समक फलन $\Id{A} \colon A \to A$, जहाँ
$\Id{A} (x) = x$ प्रत्येक $x \in A$ के लिए है, एक !!a{bijection} है। अतः
$\cardeq{A}{A}$।

\emph{सममिति।} मान लीजिए $\cardeq{A}{B}$; अर्थात कोई !!a{bijection}
$f\colon A \to B$ है। चूँकि $f$ !!{bijective} है, उसका प्रतिलोम $f^{-1}$
विद्यमान है और वह भी !!{bijective} है। अतः $f^{-1}\colon B \to A$ एक
!!a{bijection} है, इसलिए $\cardeq{B}{A}$।

\emph{संक्रामकता।} मान लीजिए कि $\cardeq{A}{B}$ और $\cardeq{B}{C}$ हैं;
अर्थात !!{bijection} $f\colon A \to B$ और $g\colon B \to
C$ हैं। तब संयोजन $\comp{f}{g}\colon A \to C$ !!{bijective} है, इसलिए
$\cardeq{A}{C}$।
\end{proof}

\begin{prop}
यदि $\cardeq{A}{B}$, तो $A$ !!{enumerable} है यदि
और केवल यदि $B$ गणनीय है।
\end{prop}

\begin{editorial}
यदि \olref[enm]{sec} सम्मिलित है, तो निम्न प्रमाण
\olref[enm]{defn:enumerable} का उपयोग करता है; अन्यथा वह
\olref[enm-alt]{defn:enumerable} का उपयोग करता है।
\end{editorial}

\begin{proof}
मान लीजिए $\cardeq{A}{B}$ है, इसलिए कोई !!{bijection} $f \colon A
\to B$ है, और मान लीजिए कि $A$ !!{enumerable} है।
\oliflabeldef{sfr:siz:enm:defn:enumerable}{
  तब या तो $A = \emptyset$ है या कोई !!a{surjective} फलन
  $g\colon \PosInt \to A$ है। यदि $A = \emptyset$, तो $B = \emptyset$
  भी है (अन्यथा कोई !!a{element}~$y \in B$ होता, लेकिन ऐसा कोई $x \in
  A$ न होता जिसके लिए $f(x) = y$)। दूसरी ओर, यदि $g\colon \PosInt \to A$
  !!{surjective} है, तो $\comp{g}{f} \colon \PosInt \to B$ भी
  !!{surjective} है। इसे देखने के लिए कोई $y \in B$ लें। चूँकि $f$
  !!{surjective} है, इसलिए कोई $x \in A$ ऐसा है कि $f(x) = y$। चूँकि
  $g$ !!{surjective} है, इसलिए कोई $n \in \PosInt$ ऐसा है कि $g(n) =
  x$। अतः
\[
(\comp{g}{f})(n) = f(g(n)) = f(x) = y
\]
और इस प्रकार $\comp{g}{f}$ !!{surjective} है। अतः $\comp{g}{f}$,
~$B$ का एक गणन है, इसलिए $B$~!!{enumerable} है।}
{
  तब या तो $A  = \emptyset$ है या कोई !!a{bijection}~$g$ है जिसका
  परिसर $A$ है और जिसका प्रांत या तो $\Nat$ है या प्राकृत संख्याओं का कोई
  आरंभिक अनुक्रम है। यदि $A = \emptyset$, तो $B =
  \emptyset$ भी है (अन्यथा कोई~$y \in B$ ऐसा होता कि ऐसा कोई $x
  \in A$ न होता जिसके लिए $f(x) = y$)। अतः मान लें कि हमारा
  !!{bijection}~$g$ है। तब $\comp{g}{f}$ एक !!a{bijection} है जिसका
  परिसर~$B$ और प्रांत वही है जो~$g$ का है (अर्थात या तो $\Nat$ या उसका
  कोई आरंभिक खंड), इसलिए $B$ !!{enumerable} है।}

यदि $B$ !!{enumerable} है, तो $A$ भी !!{enumerable} है। यह निष्कर्ष इसी
तर्क को !!{bijection} $f^{-1}\colon B \to A$ के साथ~$f$ के स्थान पर
दोहराने से मिलता है।
\end{proof}

\begin{prob}
दिखाइए कि यदि $\cardeq{A}{C}$ और $\cardeq{B}{D}$ हैं, और $A \cap B =
C \cap D = \emptyset$ है, तो $\cardeq{A \cup B}{C \cup D}$।
\end{prob}

\begin{prob}
दिखाइए कि यदि $A$ अनंत और !!{enumerable} है, तो
$\cardeq{A}{\Nat}$।
\end{prob}

\end{document}
